\documentclass{article}
\usepackage[UTF8]{ctex}
\usepackage{amsmath}
\usepackage{amssymb}
\usepackage{graphicx}
\usepackage{geometry}
\geometry{a4paper, margin=2cm}

\newcommand{\problemname}[1]{\begin{center}\Large\bfseries #1\end{center}}
\usepackage{enumitem}

\begin{document}
\problemname{数鸡兔}

\subsection*{题目描述}
蒜头君和椰菜妹一起养了一些鸡和兔子。只要给出它们的头数和脚数，请编写程序计算出鸡和兔子的数量。注意：每只鸡有 2 只脚，每只兔子有 4 只脚。

\subsection*{输入格式}
输入只有一行，包含两个用空格分开的整数，分别表示头的数量 tou 和脚的数量 jiao。这些数量都在 int 类型变量可以存储的范围内，并且保证存在整数解。

\subsection*{输出格式}
输出包括两行，每行一个整数。第一行是兔子的数量，第二行是鸡的数量。

\subsection*{样例输入}
\begin{verbatim}
88 244
\end{verbatim}

\subsection*{样例输出}
\begin{verbatim}
34
54
\end{verbatim}


\end{document}
